\documentclass[11pt,letterpaper]{article}
\usepackage[utf8]{inputenc}
\usepackage{amsmath,amssymb,amsthm}
\usepackage{geometry}
\geometry{margin=1in}

\title{APM1220: Applied Multivariate Data Analysis \\ Formative Assessment 1 (FA 1)}
\author{Atanacio, Kobe Illana A.}
\date{\today}

\begin{document}

\maketitle

\section*{Problem 2.2}
Given the matrices:
$$A = \begin{bmatrix} -1 & 3 \\ 4 & 2 \end{bmatrix}, \quad B = \begin{bmatrix} 4 & -3 \\ 1 & -2 \\ -2 & 0 \end{bmatrix}, \quad C = \begin{bmatrix} 5 \\ -4 \\ 2 \end{bmatrix}$$

\begin{enumerate}
    \item[(a)] \textbf{Find 5A:}
    Multiply each entry of matrix A by the scalar 5:
    $$5A = 5 \begin{bmatrix} -1 & 3 \\ 4 & 2 \end{bmatrix} = \begin{bmatrix} 5(-1) & 5(3) \\ 5(4) & 5(2) \end{bmatrix} = \begin{bmatrix} -5 & 15 \\ 20 & 10 \end{bmatrix}$$

    \item[(b)] \textbf{Find BA:}
    Matrix B is dimension $3 \times 2$ and matrix A is $2 \times 2$. The inner dimensions match ($2 = 2$), so the product is defined and yields a $3 \times 2$ matrix:
    $$BA = \begin{bmatrix} 4 & -3 \\ 1 & -2 \\ -2 & 0 \end{bmatrix} \begin{bmatrix} -1 & 3 \\ 4 & 2 \end{bmatrix} = \begin{bmatrix} (4)(-1) + (-3)(4) & (4)(3) + (-3)(2) \\ (1)(-1) + (-2)(4) & (1)(3) + (-2)(2) \\ (-2)(-1) + (0)(4) & (-2)(3) + (0)(2) \end{bmatrix}$$
    $$BA = \begin{bmatrix} -4 - 12 & 12 - 6 \\ -1 - 8 & 3 - 4 \\ 2 + 0 & -6 + 0 \end{bmatrix} = \begin{bmatrix} -16 & 6 \\ -9 & -1 \\ 2 & -6 \end{bmatrix}$$

    \item[(c)] \textbf{Find A'B':}
    First, find the transposes. Matrix A' is $2 \times 2$ and B' is $2 \times 3$. The inner dimensions match ($2 = 2$), yielding a $2 \times 3$ matrix:
    $$A' = \begin{bmatrix} -1 & 4 \\ 3 & 2 \end{bmatrix}, \quad B' = \begin{bmatrix} 4 & 1 & -2 \\ -3 & -2 & 0 \end{bmatrix}$$
    $$A'B' = \begin{bmatrix} -1 & 4 \\ 3 & 2 \end{bmatrix} \begin{bmatrix} 4 & 1 & -2 \\ -3 & -2 & 0 \end{bmatrix}$$
    $$A'B' = \begin{bmatrix} (-1)(4) + (4)(-3) & (-1)(1) + (4)(-2) & (-1)(-2) + (4)(0) \\ (3)(4) + (2)(-3) & (3)(1) + (2)(-2) & (3)(-2) + (2)(0) \end{bmatrix}$$
    $$A'B' = \begin{bmatrix} -4 - 12 & -1 - 8 & 2 + 0 \\ 12 - 6 & 3 - 4 & -6 + 0 \end{bmatrix} = \begin{bmatrix} -16 & -9 & 2 \\ 6 & -1 & -6 \end{bmatrix}$$

    \item[(d)] \textbf{Find C'B:}
    Matrix C' is $1 \times 3$ and B is $3 \times 2$. Inner dimensions match ($3 = 3$), yielding a $1 \times 2$ matrix:
    $$C'B = \begin{bmatrix} 5 & -4 & 2 \end{bmatrix} \begin{bmatrix} 4 & -3 \\ 1 & -2 \\ -2 & 0 \end{bmatrix} = \begin{bmatrix} (5)(4) + (-4)(1) + (2)(-2) & (5)(-3) + (-4)(-2) + (2)(0) \end{bmatrix}$$
    $$C'B = \begin{bmatrix} 20 - 4 - 4 & -15 + 8 + 0 \end{bmatrix} = \begin{bmatrix} 12 & -7 \end{bmatrix}$$

    \item[(e)] \textbf{Is AB defined?}
    Matrix A is $2 \times 2$ and matrix B is $3 \times 2$. Because the inner dimensions do not match ($2 \neq 3$), \textbf{AB is not defined}.
\end{enumerate}


\section*{Problem 2.4}
When $A^{-1}$ and $B^{-1}$ exist, prove each of the following:

\begin{enumerate}
    \item[(a)] \textbf{Prove that $(A')^{-1} = (A^{-1})'$}
    
    \textit{Proof:} By the definition of an invertible matrix, $A$ satisfies $AA^{-1} = I$ and $A^{-1}A = I$. 
    Taking the transpose of both sides of $AA^{-1} = I$ and noting that the transpose of the identity matrix is itself ($I' = I$):
    $$(AA^{-1})' = I' \implies (A^{-1})'A' = I$$
    Similarly, taking the transpose of the second relation:
    $$(A^{-1}A)' = I' \implies A'(A^{-1})' = I$$
    Since both products equal $I$, by the uniqueness of the matrix inverse, we conclude that the inverse of $A'$ must be $(A^{-1})'$. Therefore, $(A')^{-1} = (A^{-1})'$. \qed

    \item[(b)] \textbf{Prove that $(AB)^{-1} = B^{-1}A^{-1}$}
    
    \textit{Proof:} To prove this, we must show that multiplying $AB$ by $B^{-1}A^{-1}$ yields the identity matrix $I$. Consider the product $(B^{-1}A^{-1})(AB)$. Using the associative property of matrix multiplication:
    $$(B^{-1}A^{-1})(AB) = B^{-1}(A^{-1}A)B$$
    Since $A^{-1}A = I$, we can substitute $I$:
    $$B^{-1}(I)B = B^{-1}B$$
    Since $B^{-1}B = I$, we have shown that $(B^{-1}A^{-1})(AB) = I$. 
    Checking the reverse multiplication:
    $$(AB)(B^{-1}A^{-1}) = A(BB^{-1})A^{-1} = AIA^{-1} = AA^{-1} = I$$
    Thus, by definition of the inverse, $(AB)^{-1} = B^{-1}A^{-1}$. \qed
\end{enumerate}

\section*{Problem 2.8}
Given the matrix:
$$A = \begin{bmatrix} 1 & 2 \\ 2 & -2 \end{bmatrix}$$

\begin{enumerate}
    \item[(1)] \textbf{Find Eigenvalues:} 
    Solve the characteristic equation $\det(A - \lambda I) = 0$:
    $$\det \begin{bmatrix} 1 - \lambda & 2 \\ 2 & -2 - \lambda \end{bmatrix} = 0$$
    $$(1 - \lambda)(-2 - \lambda) - (2)(2) = 0$$
    $$-2 - \lambda + 2\lambda + \lambda^2 - 4 = 0 \implies \lambda^2 + \lambda - 6 = 0$$
    Factoring the quadratic yields $(\lambda + 3)(\lambda - 2) = 0$. 
    The eigenvalues are $\mathbf{\lambda_1 = 2}$ and $\mathbf{\lambda_2 = -3}$.

    \item[(2)] \textbf{Find Associated Normalized Eigenvectors:}
    \begin{itemize}
        \item For $\lambda_1 = 2$: 
        $$(A - 2I)x = \begin{bmatrix} -1 & 2 \\ 2 & -4 \end{bmatrix} \begin{bmatrix} x_1 \\ x_2 \end{bmatrix} = \begin{bmatrix} 0 \\ 0 \end{bmatrix}$$
        This gives $-x_1 + 2x_2 = 0 \implies x_1 = 2x_2$. Let $x_2=1, x_1=2$. The unnormalized vector is $[2, 1]'$.
        Length $= \sqrt{2^2 + 1^2} = \sqrt{5}$. 
        Normalized eigenvector: $\mathbf{e_1 = \begin{bmatrix} \frac{2}{\sqrt{5}} \\ \frac{1}{\sqrt{5}} \end{bmatrix}}$

        \item For $\lambda_2 = -3$: 
        $$(A - (-3)I)x = \begin{bmatrix} 4 & 2 \\ 2 & 1 \end{bmatrix} \begin{bmatrix} x_1 \\ x_2 \end{bmatrix} = \begin{bmatrix} 0 \\ 0 \end{bmatrix}$$
        This gives $2x_1 + x_2 = 0 \implies x_2 = -2x_1$. Let $x_1=1, x_2=-2$. The unnormalized vector is $[1, -2]'$.
        Length $= \sqrt{1^2 + (-2)^2} = \sqrt{5}$.
        Normalized eigenvector: $\mathbf{e_2 = \begin{bmatrix} \frac{1}{\sqrt{5}} \\ -\frac{2}{\sqrt{5}} \end{bmatrix}}$
    \end{itemize}

    \item[(3)] \textbf{Determine Spectral Decomposition:} 
    Using $A = \lambda_1 e_1 e_1' + \lambda_2 e_2 e_2'$:
    $$A = 2 \begin{bmatrix} \frac{2}{\sqrt{5}} \\ \frac{1}{\sqrt{5}} \end{bmatrix} \begin{bmatrix} \frac{2}{\sqrt{5}} & \frac{1}{\sqrt{5}} \end{bmatrix} + (-3) \begin{bmatrix} \frac{1}{\sqrt{5}} \\ -\frac{2}{\sqrt{5}} \end{bmatrix} \begin{bmatrix} \frac{1}{\sqrt{5}} & -\frac{2}{\sqrt{5}} \end{bmatrix}$$
    $$A = 2 \begin{bmatrix} \frac{4}{5} & \frac{2}{5} \\ \frac{2}{5} & \frac{1}{5} \end{bmatrix} - 3 \begin{bmatrix} \frac{1}{5} & -\frac{2}{5} \\ -\frac{2}{5} & \frac{4}{5} \end{bmatrix} = \begin{bmatrix} \frac{8}{5} & \frac{4}{5} \\ \frac{4}{5} & \frac{2}{5} \end{bmatrix} + \begin{bmatrix} -\frac{3}{5} & \frac{6}{5} \\ \frac{6}{5} & -\frac{12}{5} \end{bmatrix}$$
\end{enumerate}


\section*{Problem 2.11}
\textbf{Show that the determinant of a $p \times p$ diagonal matrix $A = \{a_{ij}\}$ with $a_{ij}=0$ for $i \neq j$, is given by $|A| = a_{11}a_{22}\cdots a_{pp}$.}

\textit{Proof:} By the definition of cofactor expansion (Definition 2A.24) along the first row, we evaluate the determinant. Since all non-diagonal elements in the first row are zero except $a_{11}$, the expansion simplifies to:
$$|A| = a_{11}A_{11} + 0 + \cdots + 0 = a_{11}A_{11}$$
where $A_{11}$ is the determinant of the $(p-1) \times (p-1)$ submatrix obtained by deleting the first row and first column of A. Because A is a diagonal matrix, the submatrix whose determinant is $A_{11}$ is also a diagonal matrix, consisting of the diagonal elements $a_{22}, a_{33}, \dots, a_{pp}$. 
Repeating this cofactor expansion process successively for each submatrix down to the final $1 \times 1$ trailing submatrix yields:
$$|A| = a_{11} \cdot a_{22} \cdot a_{33} \cdots a_{pp} = \prod_{i=1}^{p} a_{ii}$$ \qed

\section*{Problem 2.24}
Given covariance matrix $\Sigma = \begin{bmatrix} 4 & 0 & 0 \\ 0 & 9 & 0 \\ 0 & 0 & 1 \end{bmatrix}$:
\begin{enumerate}
    \item[(a)] \textbf{Find $\Sigma^{-1}$:} 
    Because $\Sigma$ is a diagonal matrix, its inverse is the diagonal matrix of the reciprocals of the original diagonal entries:
    $$\Sigma^{-1} = \begin{bmatrix} \frac{1}{4} & 0 & 0 \\ 0 & \frac{1}{9} & 0 \\ 0 & 0 & 1 \end{bmatrix}$$
    
    \item[(b)] \textbf{Eigenvalues and eigenvectors of $\Sigma$:} 
    For a diagonal matrix, the eigenvalues are the diagonal entries, and the standard basis vectors are the corresponding eigenvectors:
    $$\lambda_1 = 4, \quad e_1 = \begin{bmatrix} 1 \\ 0 \\ 0 \end{bmatrix} \quad \lambda_2 = 9, \quad e_2 = \begin{bmatrix} 0 \\ 1 \\ 0 \end{bmatrix} \quad \lambda_3 = 1, \quad e_3 = \begin{bmatrix} 0 \\ 0 \\ 1 \end{bmatrix}$$
    
    \item[(c)] \textbf{Eigenvalues and eigenvectors of $\Sigma^{-1}$:} 
    Similarly, for the diagonal inverse matrix:
    $$\lambda_1 = \frac{1}{4}, \quad e_1 = \begin{bmatrix} 1 \\ 0 \\ 0 \end{bmatrix} \quad \lambda_2 = \frac{1}{9}, \quad e_2 = \begin{bmatrix} 0 \\ 1 \\ 0 \end{bmatrix} \quad \lambda_3 = 1, \quad e_3 = \begin{bmatrix} 0 \\ 0 \\ 1 \end{bmatrix}$$
\end{enumerate}


\section*{Problem 2.30}
Given partitioned vectors and matrices:
$\mu_X = [4, 3, 2, 1]'$ partitioned into $X^{(1)}$ and $X^{(2)}$, meaning $\mu^{(1)} = \begin{bmatrix} 4 \\ 3 \end{bmatrix}$ and $\mu^{(2)} = \begin{bmatrix} 2 \\ 1 \end{bmatrix}$. 
$\Sigma_X$ is partitioned into $\Sigma_{11} = \begin{bmatrix} 3 & 0 \\ 0 & 1 \end{bmatrix}$, $\Sigma_{22} = \begin{bmatrix} 9 & -2 \\ -2 & 4 \end{bmatrix}$, and $\Sigma_{12} = \begin{bmatrix} 2 & 2 \\ 1 & 0 \end{bmatrix}$.
$A = \begin{bmatrix} 1 & 2 \end{bmatrix}$ and $B = \begin{bmatrix} 1 & -2 \\ 2 & -1 \end{bmatrix}$.

\begin{enumerate}
    \item[(a)] $\mathbf{E(X^{(1)}) = \mu^{(1)} = \begin{bmatrix} 4 \\ 3 \end{bmatrix}}$
    \item[(b)] $\mathbf{E(AX^{(1)}) = A \mu^{(1)} = \begin{bmatrix} 1 & 2 \end{bmatrix} \begin{bmatrix} 4 \\ 3 \end{bmatrix} = [1(4) + 2(3)] = [10]}$
    \item[(c)] $\mathbf{Cov(X^{(1)}) = \Sigma_{11} = \begin{bmatrix} 3 & 0 \\ 0 & 1 \end{bmatrix}}$
    \item[(d)] \textbf{Find $Cov(AX^{(1)})$:} 
    $Cov(AX^{(1)}) = A \Sigma_{11} A'$
    $$= \begin{bmatrix} 1 & 2 \end{bmatrix} \begin{bmatrix} 3 & 0 \\ 0 & 1 \end{bmatrix} \begin{bmatrix} 1 \\ 2 \end{bmatrix} = \begin{bmatrix} 3 & 2 \end{bmatrix} \begin{bmatrix} 1 \\ 2 \end{bmatrix} = [3(1) + 2(2)] = [7]$$
    \item[(e)] $\mathbf{E(X^{(2)}) = \mu^{(2)} = \begin{bmatrix} 2 \\ 1 \end{bmatrix}}$
    \item[(f)] \textbf{Find $E(BX^{(2)})$:}
    $E(BX^{(2)}) = B \mu^{(2)} = \begin{bmatrix} 1 & -2 \\ 2 & -1 \end{bmatrix} \begin{bmatrix} 2 \\ 1 \end{bmatrix} = \begin{bmatrix} 1(2) - 2(1) \\ 2(2) - 1(1) \end{bmatrix} = \mathbf{\begin{bmatrix} 0 \\ 3 \end{bmatrix}}$
    \item[(g)] $\mathbf{Cov(X^{(2)}) = \Sigma_{22} = \begin{bmatrix} 9 & -2 \\ -2 & 4 \end{bmatrix}}$
    \item[(h)] \textbf{Find $Cov(BX^{(2)})$:}
    $Cov(BX^{(2)}) = B \Sigma_{22} B' = \begin{bmatrix} 1 & -2 \\ 2 & -1 \end{bmatrix} \begin{bmatrix} 9 & -2 \\ -2 & 4 \end{bmatrix} \begin{bmatrix} 1 & 2 \\ -2 & -1 \end{bmatrix}$
    $$= \begin{bmatrix} (1)(9)+(-2)(-2) & (1)(-2)+(-2)(4) \\ (2)(9)+(-1)(-2) & (2)(-2)+(-1)(4) \end{bmatrix} \begin{bmatrix} 1 & 2 \\ -2 & -1 \end{bmatrix} = \begin{bmatrix} 13 & -10 \\ 20 & -8 \end{bmatrix} \begin{bmatrix} 1 & 2 \\ -2 & -1 \end{bmatrix}$$
    $$= \begin{bmatrix} 13(1)+(-10)(-2) & 13(2)+(-10)(-1) \\ 20(1)+(-8)(-2) & 20(2)+(-8)(-1) \end{bmatrix} = \mathbf{\begin{bmatrix} 33 & 36 \\ 36 & 48 \end{bmatrix}}$$
    \item[(i)] $\mathbf{Cov(X^{(1)}, X^{(2)}) = \Sigma_{12} = \begin{bmatrix} 2 & 2 \\ 1 & 0 \end{bmatrix}}$
    \item[(j)] \textbf{Find $Cov(AX^{(1)}, BX^{(2)})$:}
    $Cov(AX^{(1)}, BX^{(2)}) = A \Sigma_{12} B' = \begin{bmatrix} 1 & 2 \end{bmatrix} \begin{bmatrix} 2 & 2 \\ 1 & 0 \end{bmatrix} \begin{bmatrix} 1 & 2 \\ -2 & -1 \end{bmatrix}$
    $$= \begin{bmatrix} 1(2)+2(1) & 1(2)+2(0) \end{bmatrix} \begin{bmatrix} 1 & 2 \\ -2 & -1 \end{bmatrix} = \begin{bmatrix} 4 & 2 \end{bmatrix} \begin{bmatrix} 1 & 2 \\ -2 & -1 \end{bmatrix}$$
    $$= \begin{bmatrix} 4(1)+2(-2) & 4(2)+2(-1) \end{bmatrix} = \mathbf{\begin{bmatrix} 0 & 6 \end{bmatrix}}$$
\end{enumerate}

\section*{Problem 2.41}
Given $\mu_X = [3, 2, -2, 0]'$, $\Sigma_X = 3I_4$, and transformation matrix $A = \begin{bmatrix} 1 & -1 & 0 & 0 \\ 1 & 1 & -2 & 0 \\ 1 & 1 & 1 & -3 \end{bmatrix}$:
\begin{enumerate}
    \item[(a)] \textbf{Find $E(AX)$:} 
    $$E(AX) = A \mu_X = \begin{bmatrix} 1 & -1 & 0 & 0 \\ 1 & 1 & -2 & 0 \\ 1 & 1 & 1 & -3 \end{bmatrix} \begin{bmatrix} 3 \\ 2 \\ -2 \\ 0 \end{bmatrix} = \begin{bmatrix} 1(3) - 1(2) + 0 + 0 \\ 1(3) + 1(2) - 2(-2) + 0 \\ 1(3) + 1(2) + 1(-2) - 3(0) \end{bmatrix} = \mathbf{\begin{bmatrix} 1 \\ 9 \\ 3 \end{bmatrix}}$$
    
    \item[(b)] \textbf{Find $Cov(AX)$:} 
    $Cov(AX) = A \Sigma_X A' = A(3I_4)A' = 3AA'$
    First, find $AA'$ (Dimensions: $3 \times 4$ multiplied by $4 \times 3$ yields $3 \times 3$):
    $$AA' = \begin{bmatrix} 1 & -1 & 0 & 0 \\ 1 & 1 & -2 & 0 \\ 1 & 1 & 1 & -3 \end{bmatrix} \begin{bmatrix} 1 & 1 & 1 \\ -1 & 1 & 1 \\ 0 & -2 & 1 \\ 0 & 0 & -3 \end{bmatrix}$$
    $$= \begin{bmatrix} (1)(1)+(-1)(-1) & (1)(1)+(-1)(1) & (1)(1)+(-1)(1) \\ (1)(1)+(1)(-1) & (1)(1)+(1)(1)+(-2)(-2) & (1)(1)+(1)(1)+(-2)(1) \\ (1)(1)+(1)(-1) & (1)(1)+(1)(1)+(1)(-2) & (1)(1)+(1)(1)+(1)(1)+(-3)(-3) \end{bmatrix}$$
    $$= \begin{bmatrix} 2 & 0 & 0 \\ 0 & 6 & 0 \\ 0 & 0 & 12 \end{bmatrix}$$
    Multiply by scalar 3:
    $$Cov(AX) = 3 \begin{bmatrix} 2 & 0 & 0 \\ 0 & 6 & 0 \\ 0 & 0 & 12 \end{bmatrix} = \mathbf{\begin{bmatrix} 6 & 0 & 0 \\ 0 & 18 & 0 \\ 0 & 0 & 36 \end{bmatrix}}$$
    
    \item[(c)] \textbf{Which pairs of linear combinations have zero covariances?}
    Since $Cov(AX)$ is a diagonal matrix, all off-diagonal elements are 0. Therefore, \textbf{all pairs of distinct linear combinations} have zero covariances (they are mutually uncorrelated).
\end{enumerate}

\end{document}